\documentclass[utf8]{beamer}
\usepackage{amsmath}
\usepackage[ngerman]{babel}

%Zeichnen
\usepackage{tikz}
\usetikzlibrary{calc,arrows,matrix,positioning,snakes,patterns,fadings,trees,shapes,mindmap}

\title{Learn Linux the Hard Way}
\author{Christian Hundt, Moritz Schlarb}
\institute{JGU}
\usetheme{CambridgeUS}
\beamertemplatenavigationsymbolsempty

\begin{document}

\begin{frame}
\tikzfading[name=fade outside,
           left color=transparent!0,
           right color=transparent!100,
           middle color=transparent!50]
\begin{tikzpicture}[remember picture,overlay]
  \node[yshift=-3.5cm,xshift=9.5cm] at (current page.north west) {%
    \begin{tikzpicture}
        \definecolor{unired}{rgb}{0.7568, 0, 0.176}
        \draw (0.3,2) node[right]{\includegraphics[scale=0.2]{logo}};
        \draw[fill,
              pattern color=black, 
              path fading=fade outside] (4.4,2.15) rectangle (11,1.85);
        \draw (2.6,2.65) node [right] {\begin{huge}
           \textbf{\textsf{Learn Linux }}\end{huge}};
        \draw (17.3,1.4) node [left] {\textsf{Dipl.- Phys. Ch. Hundt}};
        \draw (4.7,1.25) node [right] {\begin{Huge}\textbf{\textsf{The Hard Way}}\end{Huge}};
    \end{tikzpicture}
    
   };
 \end{tikzpicture}
 
 \begin{center}
 \vspace{28pt}
 \end{center}
\begin{center}
\begin{tikzpicture}[remember picture,overlay]
  \node[yshift=-7.7cm,xshift=5.9cm] at (current page.north west) {
\includegraphics[scale=0.9]{uni.pdf}};
\end{tikzpicture}
\end{center} 
 
\end{frame}

\frame{
	\frametitle{Inhaltsverzeichnis}
	\tableofcontents
}

\section{Formales}

\begin{frame}
  \frametitle{\insertsectionhead}
  \framesubtitle{Veranstaltung}

  \begin{columns}
    \column{.66\textwidth}
      \begin{block}{}
        \begin{itemize}
          \item Veranstaltungshomepage:
            \url{https://github.com/joguinf/learnlinuxthehardway/wiki/}
          \item Mailingliste: 
            \url{learnlinuxthehardway@lists.uni-mainz.de}
          \item Archiv, RSS-Feeds, Anmeldung, Abmeldung:
            \url{https://lists.uni-mainz.de/sympa/info/learnlinuxthehardway}
        \end{itemize}
      \end{block}
    \column{.33\textwidth}
      \hfill
      \includegraphics[scale=0.2]{larry_the_cow-full.png}
  \end{columns}
\end{frame}

\begin{frame}
  \frametitle{\insertsectionhead}
  \framesubtitle{Scheinkriterien}

  \begin{block}{}
    \begin{description}
      \item[Vortrag] \hfill \\
        60 min inkl. Live-Demo der besprochenen Themen
      \item[Ausarbeitung] \hfill \\
        15 Seiten, \LaTeX{}, Vorlage wird gestellt
        \begin{itemize}
          \item \LaTeX{}-Dokumente gemeinsam erstellen:
            \url{http://writelatex.com}
        \end{itemize}
      \item[Werkstück] \hfill \\
        Virtuelle Maschine, die das fertige Ergebnis des Seminarthemas enthält
    \end{description}
  \end{block}
\end{frame}

\begin{frame}
  \frametitle{Gentoo}

  \begin{columns}
    \column{.66\textwidth}
      \begin{block}{}
        \begin{itemize}
          \item \url{http://www.gentoo.org}
          \item Source-Distribution (\textit{keine} vorkompilierten Binärpakete)
          \item Modular, Customizable, Tweakable
        \end{itemize}
      \end{block}
    \column{.33\textwidth}
      \hfill
      \includegraphics[scale=0.2]{Gentoo_logo.png}
  \end{columns}
\end{frame}

\section{Themen}

\begin{frame}
  \frametitle{\insertsectionhead}
  \framesubtitle{Basissystem}

\begin{columns}
\column{.45\textwidth}
\begin{block}{Basissystem I}
  \begin{itemize}
    \item Installationsmedium (Live-CD)
    \item Partitionierung, Wahl des Dateisystems (\texttt{ext}, \texttt{jfs}, \texttt{xfs}, \texttt{btrfs})
    \item \texttt{wget} und \texttt{tar} für das Stage3-Archiv
    \item \texttt{chroot}
  \end{itemize}
\end{block}

\column{.45\textwidth}
\begin{block}{Basissystem II}
  \begin{itemize}
    \item Portage
    \item \texttt{/etc/make.conf} und \texttt{/etc/portage/}
    \item Der Kernel - Manuell vs. Genkernel
    \vspace{3\baselineskip}
  \end{itemize}
\end{block}
\end{columns}
\end{frame}

\begin{frame}
  \frametitle{\insertsectionhead}
  \framesubtitle{Desktop}

\begin{columns}
\column{.45\textwidth}
\begin{block}{Desktop}
  \begin{itemize}
    \item Enlightenment, Xfce, Fluxbox
    \item Verschlüsselung des Heimatverzeichnisses
    \item Conky
    \item Eclipse
    \item Gaming
  \end{itemize}
\end{block}

\column{.45\textwidth}
\begin{block}{Development}
  \begin{itemize}
    \item Makefiles
    \item SWIG
    \item Ebuilds
    \item Git, SVN
    \item IRC, Mailinglisten
    \vspace{2\baselineskip}
  \end{itemize}
\end{block}
\end{columns}
\end{frame}

\begin{frame}
  \frametitle{\insertsectionhead}
  \framesubtitle{Server}

\begin{block}{Administration}
  \begin{itemize}
    \item SSH, SCP
    \item Screen
    \item Bash, Z-Shell
    \item Cron, Syslog
  \end{itemize}
\end{block}

\begin{columns}
\column{.45\textwidth}
\begin{block}{NAS}
  \begin{itemize}
    \item RAID
    \item LUKS
    \item Samba
    \item NFS
    \item rsync
  \end{itemize}
\end{block}

\column{.45\textwidth}
\begin{block}{WWW}
  \begin{itemize}
    \item LAMP (Apache, PHP, MySQL)
    \item Wordpress
    \item FTP
    \vspace{2\baselineskip}
  \end{itemize}
\end{block}
\end{columns}
\end{frame}

\begin{frame}
  \frametitle{\insertsectionhead}
  \framesubtitle{}

\begin{columns}
\column{.45\textwidth}
\begin{block}{Scientific Computing}
  \begin{itemize}
    \item MPI
    \item Cuda
    \item NumPy, SciPy
    \item GNUPlot
  \end{itemize}
\end{block}

\column{.45\textwidth}
\begin{block}{}
  \begin{itemize}
      \item
  \end{itemize}
\end{block}
\end{columns}
\end{frame}

\section{Literatur}

\begin{frame}
  \frametitle{\insertsectionhead}
  \framesubtitle{}

  \begin{block}{Gentoo}
    \begin{itemize}
      \item Gentoo Handbook - \url{http://www.gentoo.org/doc/en/handbook}
      \item Official Gentoo Wiki - \url{http://wiki.gentoo.org}
      \item Inofficial Gentoo Wiki - \url{http://en.gentoo-wiki.com}
    \end{itemize}
  \end{block}
  \begin{block}{Linux}
    \begin{itemize}
      \item Kernel Newbies - \url{http://kernelnewbies.org}
      \item Linux from Scratch - \url{http://www.linuxfromscratch.org}
      \item nixCraft: Linux Tips, Hacks, Tutorials - \url{http://www.cyberciti.biz}
    \end{itemize}
  \end{block}
\end{frame}

\end{document}
